(1) 链方程为 \(E\psi_n=-J(\psi_{n+1}+\psi_{n-1})+Vd\,\delta_{n0}\)，量子点方程 \((E-\varepsilon_d)d=V\psi_0\)。对 \(E\ne\varepsilon_d\)，消去 \(d\) 得
\[
E\psi_n=-J(\psi_{n+1}+\psi_{n-1})+U(E)\psi_0\delta_{n0},
\quad U(E)=\frac{V^2}{E-\varepsilon_d}.
\]
这是具有能量极点的Feshbach有效散射势，不是能量无关的普通 \(\delta\) 杂质；极点处应取连续散射极限。
(2) 由 \(n=\pm1\) 的离散方程知 \(\psi_0=t=1+r\)。代入 \(n=0\) 方程并用 \(E=-2J\cos k\) 化简得 \(-2iJr\sin k+U(1+r)=0\)，所以
\[
r=\frac{U}{2iJ\sin k-U},\qquad
t=\frac{2iJ\sin k}{2iJ\sin k-U}.
\]
入射和透射处群速度相同 \(v_g=2J\sin k/\hbar>0\)，故
\[
T=|t|^2=\frac{4J^2\sin^2k}{4J^2\sin^2k+U^2},\quad
R=\frac{U^2}{4J^2\sin^2k+U^2},\quad R+T=1.
\]
(3) 消去 \(U\) 后
\[
T(E)=\frac{(E-\varepsilon_d)^2}{(E-\varepsilon_d)^2+
\Gamma(E)^2},\quad
\Gamma(E)=\frac{V^2}{2J\sin k(E)}.
\]
若点能在带内且 \(V>0\)，在 \(E=\varepsilon_d\) 有 \(T=0,R=1\)，反射由侧点的破坏性干涉造成。窄共振时可用 \(\Gamma(\varepsilon_d)\) 作局部半宽参数。先取 \(V\to0\) 且 \(E\ne\varepsilon_d\) 得 \(T\to1\)，先令 \(E=\varepsilon_d\) 再取 \(V\to0^+\) 得 \(T=0\)；在严格 \(V=0\) 时量子点脱耦，恢复 \(T=1\)。
(4) 对 \(|E|>2J\)，无限链局域Green函数为 \(g_{00}(E)=\operatorname{sgn}(E)/\sqrt{E^2-4J^2}\)，点态自洽条件
\[
E-\varepsilon_d=V^2\frac{\operatorname{sgn}(E)}{\sqrt{E^2-4J^2}}.
\]
在 \(E>2J\) 左侧随 \(E\) 严格递增，右侧由正无穷降至零，恰有一个正根；在 \(E<-2J\) 相反亦恰有一个负根。对 \(\varepsilon_d=0\)，平方并保留带外条件得
\[
E_\pm=\pm\sqrt{2J^2+\sqrt{4J^4+V^4}}.
\]
当 \(V\to0^+\)，两支分别趋近 \(\pm2J\) 的带边；对应态的链上波函数指数衰减，满足正规化而非散射边界。
